% Контрольна робота з прикладної криптології — варіант
% Компіляція: xelatex / lualatex (fontspec) або pdflatex (T2A).
% Відповіді: розкоментуйте \answerstrue нижче.
\documentclass[12pt,a4paper]{article}
\usepackage{yfonts}
\usepackage{emoji}

\usepackage{iftex}
\ifPDFTeX
  \usepackage[T2A]{fontenc}
  \usepackage[utf8]{inputenc}
  \usepackage[ukrainian]{babel}
\else
  \usepackage{fontspec}
  \IfFontExistsTF{Times New Roman}
    {\setmainfont{Times New Roman}}
    {\setmainfont{Liberation Serif}}
  \usepackage{polyglossia}
  \setdefaultlanguage{ukrainian}
\fi

\usepackage{amsmath,amssymb}
\usepackage[margin=2cm]{geometry}
\usepackage{array}
\usepackage{enumitem}

\newif\ifanswers
% \answerstrue   % <- розкоментувати, щоб отримати версію з відповідями

\newcommand{\variant}{1}
\newcommand{\F}{\mathbb{F}}
\usepackage{environ}
\NewEnviron{answer}{\ifanswers\par\smallskip\noindent\textit{Відповідь.}\ \BODY\par\fi}

\setlength{\parindent}{0pt}
\pagestyle{empty}

\begin{document}

\begin{center}
  {\Huge\textfrak{Cryptography CW }}   % Gothisch — класична «текстура»
\end{center}
\begin{center}
  \textswab{ Variant \emoji{new-moon-face}}
  \ifanswers\\[2pt]\textit{(з відповідями)}\fi
\end{center}

\vspace{4pt}

% ---------------------------------------------------------------- 1
\textbf{1.} Відомо, що слово ОКАРУН було зашифроване афінним шифром у слово
АЙЗИХШ. Розшифруйте шифртекст ЙЗЖБМЯ.

\smallskip
{\small
\setlength{\tabcolsep}{1.7pt}
\begin{tabular}{|*{33}{c|}}
\hline
А&Б&В&Г&Ґ&Д&Е&Є&Ж&З&И&І&Ї&Й&К&Л&М&Н&О&П&Р&С&Т&У&Ф&Х&Ц&Ч&Ш&Щ&Ь&Ю&Я\\ \hline
0&1&2&3&4&5&6&7&8&9&10&11&12&13&14&15&16&17&18&19&20&21&22&23&24&25&26&27&28&29&30&31&32\\ \hline
\end{tabular}}

\begin{answer}
$E(x)=ax+b \bmod 33$. З пар О$\to$А ($18\to0$), К$\to$Й ($14\to13$):
$4a\equiv-13\equiv20 \pmod{33}$, $4^{-1}=25$, отже $a=20\cdot25\bmod33=5$,
$b=-18\cdot5\bmod 33=9$. Ключ $(a,b)=(5,9)$ узгоджується з усіма шістьма парами.
$a^{-1}=20$, $D(y)=20(y-9)\bmod 33$.
ЙЗЖБМЯ $=(13,9,8,1,16,32)\mapsto(14,0,13,5,8,31)=$ \textbf{КАЙДЖЮ}.
\end{answer}

\medskip
% ---------------------------------------------------------------- 2
\textbf{2.} Хлопець Нацуме на святі духів шукає серед йокаїв того, хто був би
добрим і водночас знав його бабусю Рейко. Відомо, що серед усіх йокаїв, які
прийшли на свято, добрий кожен третій, а Рейко знав кожен четвертий. Проте серед
тих, хто знав Рейко, добрим є лише кожен п'ятий. Яка невизначеність щодо того,
що перший зустрічний йокай відповідає сподіванням Нацуме, $H(X,Y)$
($X$~--- доброта йокая, $Y$~--- знайомство з Рейко)? Яка невизначеність щодо
доброти йокая залишиться, якщо Нацуме дізнається, що той знав Рейко,
$H(X\,|\,Y=\text{знав Рейко})$? Скільки інформації про доброту йокая Нацуме при
цьому отримає?

\begin{answer}
$P(Y{=}1)=\tfrac14$, $P(X{=}1)=\tfrac13$, $P(X{=}1\,|\,Y{=}1)=\tfrac15$.
Спільний розподіл: $P(1,1)=\tfrac1{20}$, $P(1,0)=\tfrac13-\tfrac1{20}=\tfrac{17}{60}$,
$P(0,1)=\tfrac14-\tfrac1{20}=\tfrac15$, $P(0,0)=\tfrac{7}{15}$.
$H(X,Y)\approx\mathbf{1{,}709}$ біт.
$H(X\,|\,Y{=}1)=h(\tfrac15)\approx\mathbf{0{,}722}$ біт;
$H(X)=h(\tfrac13)\approx0{,}918$ біт, отже отримана інформація
$H(X)-H(X\,|\,Y{=}1)\approx\mathbf{0{,}196}$ біт.
\end{answer}

\medskip
% ---------------------------------------------------------------- 3
\textbf{3.} Обчислити стан ЛРЗ над полем $\F_7$ у 10-му такті, якщо початкове
заповнення та поліном зворотного зв'язку цього ЛРЗ мають відповідно вид
$(1,0,2,3)$, \ $f(x)=x^4+x^3+3$ (нумерація тактів роботи ЛРЗ починається з
числа~0). Обчислити також максимальний можливий період ЛРЗ із поліномом
зворотного зв'язку степеня~4 над $\F_7$. Відповідь обґрунтувати.

\begin{answer}
$f$ розглядається як характеристичний поліном:
$s_{t+4}=-(s_{t+3}+3s_t)=6s_{t+3}+4s_t \bmod 7$, стан у такті $t$~---
$(s_t,s_{t+1},s_{t+2},s_{t+3})$.
Послідовність: $1,0,2,3,\,1,6,2,3,1,2,6,6,5,3$.
Стан у такті 10: $\mathbf{(6,6,5,3)}$.
Максимальний період $7^4-1=\mathbf{2400}$; досягається тоді й лише тоді, коли
поліном примітивний (таких $\varphi(2400)/4=160$).
Даний $f$ незвідний (не має коренів і квадратичних множників над $\F_7$), але не
примітивний: $\operatorname{ord} f=480$, тож період цієї послідовності~--- 480.
\end{answer}

\medskip
% ---------------------------------------------------------------- 4
\textbf{4.} Гіпотетичний $S$-блок

\smallskip
{\setlength{\tabcolsep}{4pt}
\begin{tabular}{|*{16}{c|}}
\hline
12&0&8&4&6&15&11&1&10&9&3&14&5&2&12&7\\ \hline
1&15&13&8&10&3&5&14&12&7&6&11&0&4&9&2\\ \hline
7&11&4&1&9&12&14&0&2&6&10&13&15&3&5&8\\ \hline
2&14&1&7&4&10&8&13&15&12&9&0&3&5&6&11\\ \hline
\end{tabular}}

\smallskip
використовується аналогічно стандартним $S$-блокам алгоритму DES, перетворюючи
вхідний 6-бітовий вектор у вихідний 4-бітовий. Встановіть відповідність між
вхідними векторами 100110, 011011, 001100, 110101 та виходами 1001, 0001, 0100,
1011.

\begin{answer}
Рядок~--- біти $b_1b_6$, стовпчик~--- $b_2b_3b_4b_5$ (нумерація з 0).\\
$100110$: рядок 2, стовпчик 3 $\to 1=\mathbf{0001}$;\quad
$011011$: рядок 1, стовпчик 13 $\to 4=\mathbf{0100}$;\\
$001100$: рядок 0, стовпчик 6 $\to 11=\mathbf{1011}$;\quad
$110101$: рядок 3, стовпчик 10 $\to 9=\mathbf{1001}$.
\end{answer}

\end{document}